\documentclass[11pt,a4paper]{article}

\usepackage[margin=1in]{geometry}
\usepackage{amsmath,amssymb,amsfonts,mathtools,bm}
\usepackage{booktabs,tabularx,array}
\usepackage{enumitem}
\usepackage{hyperref}
\usepackage[dvipsnames]{xcolor} % Enables extra named colors (single load: fixes v2 option clash)
\usepackage{microtype}
\usepackage{graphicx}
%\usepackage{xcolor}
\usepackage{soul}

\hypersetup{
  colorlinks=true,
  linkcolor=blue,
  citecolor=blue,
  urlcolor=blue
}

\newcommand{\R}{\mathbb{R}}
\newcommand{\SE}{\mathrm{SE}(3)}
\newcommand{\SO}{\mathrm{SO}(3)}
\newcommand{\softplus}{\operatorname{softplus}}
\newcommand{\sigmoid}{\operatorname{sigmoid}}
\newcommand{\normalize}{\operatorname{normalize}}
\newcommand{\stopgrad}{\operatorname{sg}}
\newcommand{\clip}{\operatorname{clip}}
\newcommand{\Sshade}{\mathcal{S}}
\newcommand{\Lshade}{\mathcal{L}_{\mathrm{shade}}}
\newcommand{\Lgeo}{\mathcal{L}_{\mathrm{geo}}}
\newcommand{\LGA}{\mathcal{L}_{\mathrm{GA}}}
\newcommand{\Loss}{\mathcal{L}}

\title{\textbf{LA-DUSt3R: Light-Aware, Pose-Free Multi-View 3D Reconstruction}\\
\large A Physically Informed Extension of DUSt3R (revised v4.1)}
\author{}
\date{}

\begin{document}
\maketitle

\begin{abstract}
We propose LA-DUSt3R, a light-aware extension of DUSt3R for feed-forward
multi-view 3D reconstruction under controlled near-field illumination.
The method combines DUSt3R's pose-free pairwise pointmap prediction and
geometric global alignment with light-conditioned features, pointmap-derived
surface normals, diffuse-albedo prediction, a differentiable direct-illumination
image-formation model, and photometric supervision during network training.
The initial formulation uses camera-attached point lights with recorded
renderer-space positions and controlled intensity settings, a fixed and
consistent rendering pipeline, and approximately Lambertian materials.
These assumptions define a controlled synthetic experiment, not a
radiometrically calibrated capture system. More general reflectance and
light-transport models are deferred.

A central design requirement is that the training-time photometric objective
needs no camera poses. For each ordered image pair, DUSt3R's reference head
predicts a pointmap in that view's own camera frame; running the swapped pair
yields the other view's camera-local pointmap. Shading is evaluated on these
camera-local pointmaps against camera-local light coordinates, so no relative
pose enters the photometric loss. Ground-truth poses are used only to build
the geometric supervision targets already required by DUSt3R.

Photometric supervision is used to shape the network's learned geometric
representation during training. At inference, the trained network predicts
pairwise pointmaps and geometric confidence, and the standard DUSt3R geometric
global-alignment procedure estimates camera poses, focal lengths, and globally
consistent geometry. No photometric-family objective, albedo optimization,
shading refinement, or depth-residual optimization is added to scene
optimization. Albedo and photometric-reliability outputs remain auxiliary
predictions and are not required by geometric global alignment. Metric scale
is supplied by metric geometric supervision; recovery of absolute scale from
illumination is not investigated.
\end{abstract}

\tableofcontents

\section{Motivation and Research Objective}

DUSt3R formulates unconstrained multi-view reconstruction as pairwise
regression of dense pointmaps, without requiring camera poses or intrinsics
as inputs. For image collections, its global alignment procedure reconciles
pairwise predictions in a common reference frame
\cite{wang2024dust3r}. This provides a strong geometric prior, but the
standard reconstruction objective is primarily geometric and does not
explicitly model the illumination process that generated the input images.

Recent point-based inverse-rendering methods demonstrate that known
illumination can provide useful constraints on geometry and appearance.
Gjerde et al.\ formulate differentiable point rendering under known
illumination, while Jiang et al.\ develop differentiable light transport
with Gaussian surfels using an adapted radiosity formulation
\cite{gjerde2026point,jiang2025radiositygs}. LA-DUSt3R investigates a
different operating point: retain DUSt3R's feed-forward, pose-free geometric
prior and use a comparatively simple direct-illumination model as an
additional training-time supervision signal.

The central research question is:

\begin{quote}
Can controlled-illumination supervision during training improve the
geometric reconstruction produced by a pose-free, feed-forward multi-view
model, when inference retains standard geometry-driven global alignment?
\end{quote}

The initial contribution is not a general inverse-rendering system. It is a
controlled integration of a learned geometric prior and a differentiable
per-pixel image-formation model. The first version deliberately assumes
Lambertian reflectance and direct illumination. Specular BRDFs, indirect
illumination, radiosity, and more general light transport are outside the
initial scope.

\subsection{Research hypotheses}

The design is organized around the following falsifiable hypotheses:
\begin{enumerate}[label=\textbf{H\arabic*.}]
  \item Photometric supervision under controlled near-field illumination
  improves surface orientation and fine-scale geometry beyond geometric
  DUSt3R supervision alone.
  \item Jointly learning diffuse albedo and photometric reliability with
  pointmap-derived shading supervision improves the geometric representation,
  while surface orientation remains derived from the pointmap. The initial
  study evaluates per-view albedo; cross-view albedo tying is deferred.
  \item Light-conditioned feature representations improve reconstruction
  \emph{when the camera-frame illumination varies across views} (for
  example, randomized light offsets and intensities in synthetic data),
  beyond the effect of additional synthetic training data alone. For a rigid
  rig with constant light offsets this hypothesis is not testable, since the
  conditioning signal is constant.
\end{enumerate}

\section{Problem Formulation and Assumptions}

\subsection{Inputs and outputs}

Let a scene be observed by $N$ cameras. The input consists of images
\[
  \mathcal{I}=\{I_n\}_{n=1}^{N},
  \qquad
  I_n:\Omega_n\rightarrow\R^3,
\]
and recorded per-image light sets
\[
  \mathcal{Q}_n
  =
  \left\{
    (\bm{\ell}_{n,k},\bm{E}_{n,k})
  \right\}_{k=1}^{K_n}.
\]
Here, $\bm{\ell}_{n,k}\in\R^3$ is the position of point light $k$ in
camera $n$'s coordinate frame, and
$\bm{E}_{n,k}\in\R_{\geq 0}^{3}$ is its RGB renderer-intensity parameter under the selected
rendering configuration; it is not assumed to be a physically calibrated
radiant intensity.

The camera poses and intrinsics are not supplied at inference. The network
predicts pairwise pointmaps and geometric confidence, together with auxiliary
diffuse-albedo and photometric-reliability outputs. Surface normals are
derived from pointmaps when needed for training-time shading supervision or
evaluation. Inference-time geometric global alignment uses the pointmaps and
geometric confidence; it does not optimize albedo, reliability, or shading.

The initial setting assumes:
\begin{itemize}
  \item camera-attached point lights (e.g.\ a simulated flash or a
  camera-mounted light rig), with their positions and renderer settings recorded;
  \item all $K_n$ lights of a view are on simultaneously, so each image is a
  single summed-illumination observation (one-light-at-a-time multi-exposure
  capture is not assumed);
  \item a consistent rendering and image-encoding pipeline, with fixed exposure
  and tone-mapping settings; if the renderer provides linear RGB, it is used
  directly, otherwise a fixed documented conversion is applied;
  \item direct-dominant illumination and diffuse-dominant materials; ambient
  illumination is disabled or held fixed and recorded;
  \item known foreground masks or a reliable procedure for obtaining them;
  \item metric camera-relative light positions and controlled renderer-intensity
  parameters (scene geometry and light positions are specified in meters).
  These light parameters are supplied to the light-conditioned network at
  inference when that conditioning pathway is enabled; they are not used in
  a post-alignment photometric objective.
\end{itemize}
These assumptions define the controlled Level-2 synthetic setting. Non-Lambertian
materials and indirect illumination are acknowledged limitations and are
reserved for later extensions.

\subsection{Training-time and inference-time information}

Ground-truth relative poses are available in the controlled training data
and are used only for the geometric supervision that DUSt3R already
requires (ground-truth pointmaps). They are not network inputs, and the
photometric losses do not use them.

\begin{center}
\begin{tabularx}{\textwidth}{>{\raggedright\arraybackslash}p{0.31\textwidth}XX}
\toprule
Quantity & Training & Inference\\
\midrule
Images and recorded light parameters & Available & Available when light conditioning is enabled\\
Ground-truth relative poses & Used only to build GT pointmaps & Unavailable\\
Ground-truth depth/geometry & Available & Unavailable\\
Ground-truth albedo & Available for synthetic supervision & Unavailable\\
Predicted pointmaps and geometric confidence & Available & Used by geometric global alignment\\
Predicted albedo and reliability & Supervised auxiliary outputs & Optional outputs; not used by global alignment\\
Camera-local pointmaps and derived normals & Used by shading and normal losses & Not required for scene optimization\\
Differentiable shading objective & Training only & Not evaluated\\
\bottomrule
\end{tabularx}
\end{center}

The photometric objective is pose-free during training. At inference, the
reconstruction pipeline is geometry-driven: no shading loss or other
photometric-family constraint is evaluated during scene optimization.

No camera poses are supplied as input at inference; they are estimated by
geometric global alignment. The photometric loss is also pose-free during
training, because it uses camera-local pointmaps and camera-local light
coordinates rather than pose-based frame conversion.

\subsection{What shading can and cannot identify}
\label{sec:ident}

This subsection fixes expectations for the hypotheses above.
\begin{itemize}
  \item \textbf{Single effective light, per-view albedo.} With known
  albedo, one observation gives one scalar equation per color channel in
  the unit normal, which constrains the normal only to a cone around the
  light direction, not to a unique vector. Pointmap derivatives supply the
  tangent-plane structure, so shading acts as a shape-from-shading
  constraint coupled to the geometric loss. It is not a free normal
  measurement.
  \item \textbf{Simultaneous lights.} $K$ lights that are on together are
  summed into one image, so they give no more equations per pixel than one
  effective light. Therefore, increasing $K$ is not expected to behave
  like adding photometric-stereo exposures. Genuine multi-light constraints
  require one-light-at-a-time exposures per viewpoint, which v1 does not
  assume.
  \item \textbf{Potential cross-view constraint.} With camera-attached
  lights, the same surface point can be observed under different
  camera-relative illumination directions. If albedo is tied across views,
  these observations could provide a multi-view photometric-stereo constraint.
  Cross-view albedo tying is deferred because reliable correspondences and
  albedo consistency introduce additional variables; it is not needed by the
  initial training objective or inference pipeline.
  \item \textbf{Scale is not a target.} Under an inverse-square approximation, scaling all depths and
  light-relative distances by $s$ while scaling albedo by $s^2$ can leave
  image intensity approximately unchanged. This is an illustrative
  ambiguity, not an assumption about the Level-2 renderer; the actual
  relationship depends on its attenuation and appearance settings.
  Metric scale from the Level-2 photometric signal is not treated as
  identifiable or recoverable in this study. We do not investigate it: scale comes from the metric
  ground-truth geometric supervision. No photometric term is evaluated
  during inference-time global alignment.
\end{itemize}

\section{Architecture Overview}

LA-DUSt3R consists of three functional components:
\begin{enumerate}
  \item \textbf{Light-conditioned DUSt3R backbone:} a DUSt3R encoder and
  cross-view decoder augmented with per-view illumination conditioning.
  \item \textbf{Dense prediction:} DUSt3R pointmaps and confidence are
  retained, while diffuse-albedo and photometric-reliability heads are added. Normals are derived from the
  pointmaps rather than predicted.
  \item \textbf{Training-time differentiable geometry and shading operators:}
  camera-local pointmaps are converted to normals, and a direct-illumination
  Lambertian-like model supplies photometric supervision during training.
  These operators are not part of inference-time scene optimization.
\end{enumerate}

\subsection{Light-conditioned feature extraction}

For each view, the recorded light set is encoded independently. For
light $k$ in view $n$, define
\[
  \bm{\tau}_{n,k}
  =
  \operatorname{MLP}_{\phi}
  \left(
    \left[
      \gamma(\bm{\ell}_{n,k}),
      \log(\bm{1}+\bm{E}_{n,k})
    \right]
  \right)
  \in\R^{d},
\]
where $\gamma(\cdot)$ denotes a positional/Fourier feature mapping and
$d$ is the decoder feature dimension.

\paragraph{Default: light tokens in self-attention only.}
The light tokens of view $n$ are appended to view $n$'s token sequence
(without positional embedding) and are visible as keys \emph{only} to view
$n$'s own self-attention. They are excluded from the keys of cross-attention,
so view 2's tokens never attend to view 1's camera-frame light coordinates.
Unlike a global modulation vector, tokens let each image token use the light
information in a spatially varying way. This matters because a pixel's
light direction depends on where it sits relative to the light. Note that
relating $\bm\ell$ to pixel rays requires the focal length, which the network
must infer implicitly.

\paragraph{Ablation: FiLM.}
View-local feature modulation
$F_n'=\operatorname{FiLM}(F_n,e(\mathcal{Q}_n))$, with $e(\mathcal{Q}_n)$
an aggregate of the encoded lights, is retained as an ablation. A single
global modulation vector cannot encode per-pixel light geometry, so we
expect it to be weaker.

Implementation detail: at every conditioned self-attention block, append
the view-local light tokens to that view's image-token sequence; image tokens
may attend to these light tokens, while light tokens are excluded from all
cross-view attention key/value sets. Remove light tokens before the dense
prediction head so the spatial token count and patch ordering remain exactly
those expected by the DUSt3R head. A learned null-light embedding and light
dropout are used when training with mixed illuminated and non-illuminated
data. Encoder-level versus
decoder-level conditioning is an experimental variable, not a fixed
assumption.

\paragraph{When conditioning is informative.}
For a rigid flash or ring-light rig, $\bm\ell_{n,k}$ is the same in every
view, so the conditioning signal is constant and may be mostly absorbed into
the weights. Hypothesis H3 is therefore tested on synthetic data with
randomized light offsets and intensities (reported as such), and separately
on a fixed-rig configuration to measure whether shading supervision helps
even when explicit light conditioning carries little view-dependent information.

\subsection{Dense prediction heads}

For an ordered image pair $(I_a,I_b)$, DUSt3R predicts pointmaps for both views in
the camera frame of the first image:
\[
  X^{a,a},\ X^{b,a}:\Omega\rightarrow\R^3,
\]
where $X^{a,a}$ is the reference view's pointmap in its own camera frame.
For brevity, when the pair is $(I_1,I_2)$ we write $X^{v,1}$, $v\in\{1,2\}$.
The extended prediction head produces
\[
  \left(
    \widetilde X^{v,1},
    \widetilde C^{v,1},
    a_\rho^v,
    \widetilde W^v
  \right)
  =
  H_\theta(G^v),
\]
where $G^v$ denotes the decoder representation for view $v$. For lit
metric samples the pointmap output is interpreted directly in meters,
$X^{v,1}=\widetilde X^{v,1}$. The per-pixel parameterizations are:
\begin{align}
  C^{v,1} &= 1+\exp(\widetilde C^{v,1}),\\
  \rho^v &= \rho_{\max}\sigmoid(a_\rho^v),\\
  W^v &= \sigmoid(\widetilde W^v).
\end{align}
Here $C$ is the DUSt3R geometric confidence, $\rho$ is bounded RGB
diffuse albedo in $[0,\rho_{\max}]$ (with $\rho_{\max}=1$ under
the normalized reflectance convention), and $W\in(0,1)$ is a
photometric-reliability weight. The bound is a modeling convention for
normalized diffuse reflectance; renderer albedo metadata must use the same
convention.
The pointmap, confidence, albedo, and reliability outputs require
$3+1+3+1=8$ channels per pixel.
The original pointmap and confidence
parameters are initialized from a DUSt3R checkpoint; the newly introduced
albedo and reliability outputs are initialized separately.

\section{Metric Data Convention and Camera-Local Pointmaps}

\subsection{Metric data convention (scale is not a research target)}

DUSt3R's original training mixture does not provide uniform metric ground
truth and therefore uses scale-normalized geometric supervision. In the
controlled synthetic subset, scene geometry, camera poses, and light
positions are specified in meters. For these lit samples, the pointmap
regression target is metric:
\[
  \ell_{\mathrm{3D}}(X^{v,1},\bar X^{v,1}),
  \qquad \bar X^{v,1}\text{ expressed in meters}.
\]
Thus, metric scale is supplied by geometric supervision. We do not train an
additional pair-scale head, introduce a new global metric-scale variable,
or attempt to recover scale from illumination. Null-light and legacy
samples retain the original normalized DUSt3R objective.

The original pretrained DUSt3R baseline may remain scale-ambiguous. For
shape comparisons, any evaluation-time scale alignment applied to that
baseline must be reported explicitly; absolute-scale metrics should also be
reported without that correction. No oracle scale is supplied to the
LA-DUSt3R network or its inference-time optimizer. Metric-supervised
LA-DUSt3R is trained to predict pointmaps in meters on the controlled
synthetic subset.

\subsection{Camera-local pointmaps for training-time shading supervision}

The shading model needs each view's geometry in that view's own camera
frame, where its light set $\mathcal{Q}_n$ is defined. DUSt3R already provides
this without any pose. Run the network on both orderings of a pair:
\[
  \boxed{
  X^{1,1}=\text{reference pointmap of }F(I_1,I_2),
  \qquad
  X^{2,2}=\text{reference pointmap of }F(I_2,I_1).
  }
\]
DUSt3R's training already feeds each pair together with its inversion, and
global alignment already uses both orderings, so this adds no new data
requirement. Each reference pointmap is camera-local by construction.
During training, shading is therefore evaluated per view:
\[
  \widehat I_1
  =
  \Sshade(X^{1,1},n_X(X^{1,1}),\rho^1,\mathcal{Q}_1),
  \qquad
  \widehat I_2
  =
  \Sshade(X^{2,2},n_X(X^{2,2}),\rho^2,\mathcal{Q}_2).
\]
Normals are computed from the camera-local pointmap, so pointmap derivatives
and light coordinates share one coordinate frame. No camera-local light
coordinate is compared with a point expressed in another camera's frame.

\section{Training-Time Differentiable Geometry and Image Formation}

\subsection{Depth back-projection}

For a camera with centered principal point and focal length $f_n$, define
\[
  K_n =
  \begin{bmatrix}
    f_n & 0 & W_n/2\\
    0 & f_n & H_n/2\\
    0 & 0 & 1
  \end{bmatrix}.
\]
Given depth $D_n(\bm p)$ at pixel
$\bm p=(u,v)^\top$, the camera-local 3D point is
\[
  \bm x_n(\bm p)
  =
  D_n(\bm p)
  K_n^{-1}
  \begin{bmatrix}u\\v\\1\end{bmatrix}.
\]
When the network pointmap is already available, it is used directly. This
back-projection relation is included to make the connection between depth,
camera-local points, and pointmap-derived normals explicit; it does not
introduce a separate inference-time depth-refinement variable.

\subsection{Geometric normals}

Normals are derived from the camera-local pointmap using finite
differences:
\[
  \bm n_X(\bm p)
  =
  \mathcal N(\bm x)(\bm p)
  =
  \normalize
  \left(
    \frac{\partial\bm x}{\partial u}
    \times
    \frac{\partial\bm x}{\partial v}
  \right).
\]
Here $\mathcal N$ denotes the pointmap-to-normal operator. The sign is chosen to satisfy the camera-facing convention,
\[
  \bm n_X(\bm p)^\top
  \frac{\bm c_n-\bm x_n(\bm p)}
       {\|\bm c_n-\bm x_n(\bm p)\|}
  >0,
\]
where $\bm c_n$ is the camera center in its own frame (the origin).

\paragraph{Discontinuity mask.}
Pixels near depth discontinuities are excluded from normal-based losses. A
criterion of the form $\|\nabla D\|/D<\tau$ also rejects legitimate
grazing-angle surfaces, which are among the most informative for shading.
Instead, compare the 3D spacing of neighboring points with the expected
pixel footprint $D/f$:
\[
  M_{\mathrm{edge}}(\bm p)
  =
  \mathbb{1}
  \left[
    \max_{\bm q\in\mathcal N_4(\bm p)}
    \left\|\bm x_n(\bm q)-\bm x_n(\bm p)\right\|
    <
    \tau_e\,\frac{D_n(\bm p)}{f_n}
  \right].
\]
A surface at incidence angle $\theta$ has spacing $\approx D/(f\cos\theta)$,
so $\tau_e\approx4$ keeps surfaces up to roughly $75^\circ$ from the
view direction while still rejecting depth jumps.

A lightly smoothed depth map may be used for finite differences, while
the unsmoothed geometry remains available to the geometric objective.
Finite differences must be validity-aware: neighborhoods
crossing invalid pixels, foreground boundaries, or discontinuities
are excluded from the shading objective. All binary validity and
discontinuity masks are computed under stop-gradient (detached) semantics;
predicted geometry must not reduce its loss by manipulating the mask. Central differences are used
in smooth interior regions, with one-sided differences only where their
validity can be established.

\subsection{Direct-illumination Lambertian model}
\label{sec:lambert}

For a camera-local surface point $\bm x$, unit normal $\bm n$, RGB albedo
$\bm\rho$, and point light $k$, define
\[
  \bm d_k=\bm\ell_k-\bm x,\qquad
  r_k=\|\bm d_k\|_2,\qquad
  \bm\omega_k=\frac{\bm d_k}{r_k}.
\]
The initial direct-illumination model is
\[
  \boxed{
  \Sshade(\bm x,\bm n,\bm\rho,\mathcal Q,V)
  =
  \frac{g\bm\rho}{\pi}
  \odot
  \sum_{k=1}^{K}
  \bm E_k V_k\,\mathcal A_k^{\mathrm{rend}}(r_k)
  \max(0,\bm n^\top\bm\omega_k).
  }
\]
Here $g$ is a fixed global rendering gain, $V_k\in[0,1]$ is visibility,
and $\odot$ denotes channel-wise multiplication. The attenuation
$\mathcal A_k^{\mathrm{rend}}$ is the configured renderer-space
attenuation function (or a documented approximation to it); it is not
assumed to be an inverse-square law. In the first controlled experiment,
visibility is set to $V_k=1$, and exposure, tone mapping, and renderer
settings are held fixed across the dataset.
The analytic shading surrogate is differentiable with respect to point
positions, normals, and albedo. It approximates the selected renderer
configuration; it is not claimed to reproduce the complete rendering
pipeline.

This model is intentionally Lambertian and direct-only. The use of
robust losses and reliability weights is intended to reduce the influence
of pixels that violate the simplified model; it does not make the model
physically complete. Specular reflection, indirect illumination,
translucency, and radiosity are future extensions.

\paragraph{Real captures.}
Real captures contain ambient illumination and additional camera-response
uncertainties. Handling these effects is deferred; possible future options
include an ambient term or flash/no-flash differencing, but neither is part
of the Level-2 synthetic study.
The initial synthetic subset disables ambient illumination where supported.
If the renderer retains a nonzero ambient term, its setting is fixed and
recorded; the analytic model is then treated as an approximation, and
strongly affected samples are excluded from the primary experiment.

\section{Training Objective}

Let $\Omega_v$ denote valid foreground pixels for view $v$. Training runs on
ordered pairs; for each ordered pair $(I_a,I_b)$, the photometric and
derived-normal losses use the \emph{reference} outputs ($v=a$), which are
camera-local. Both orderings of every pair are used. The initial training objective is
\[
  \boxed{
  \Loss_{\mathrm{train}}
  =
  \lambda_{\mathrm{geo}}\Loss_{\mathrm{geo}}
  +\lambda_{\mathrm{nrm}}\Loss_{\mathrm{nrm}}
  +\lambda_{\mathrm{shade}}\Loss_{\mathrm{shade}}
  +\lambda_{\rho}\Loss_{\rho}.
  }
\]
There is no separate normal head and no normal-consistency term between two
normal sources. Cross-view albedo consistency is outside the initial
training objective. The photometric term is initially small and is warmed up only after the
geometric prediction becomes stable.
Operationally, $\lambda_{\mathrm{shade}}$ is held at zero until the
validation $\Loss_{\mathrm{geo}}$ changes by less than $1\%$ over an
epoch (or a fixed step count $T_0$ for reproducibility), then ramped
linearly.

\subsection{Confidence-weighted geometric regression}

The geometric loss follows the confidence-weighted pointmap regression
principle of DUSt3R:
\[
  \Loss_{\mathrm{geo}}
  =
  \sum_{v\in\{1,2\}}\sum_{\bm p\in\Omega_v}
  \left[
    C^{v,1}(\bm p)\,
    \ell_{\mathrm{3D}}
    \left(
      X^{v,1}(\bm p),
      \bar X^{v,1}(\bm p)
    \right)
    -\alpha_C\log C^{v,1}(\bm p)
  \right].
\]
The precise regression norm and normalization protocol should match the
selected DUSt3R training configuration. For lit metric samples the regression is on metric pointmaps with no scale
normalization; for null-light or legacy samples, the
original normalized geometric objective is retained.

\subsection{Derived-normal supervision}

Pointwise 3D regression tolerates high-frequency errors that strongly
corrupt finite-difference normals, yet shading depends on exactly those
derivatives. To supervise them without a normal head or ground-truth normals,
apply the same operator $\mathcal N$ (including the same smoothing and masks)
to the reference output and to the ground-truth pointmap:
\[
  \Loss_{\mathrm{nrm}}
  =
  \sum_{\bm p\in\Omega_a}
  M_{\mathrm{edge}}^{\mathrm{gt}}(\bm p)\,M_{\mathrm{edge}}(\bm p)
  \left(
    1-\bm n_X\!\left(X^{a,a}\right)(\bm p)^{\!\top}
    \bm n_X\!\left(\bar X^{a,a}\right)(\bm p)
  \right).
\]
This is a gradient-matching term in the spirit of scale-shift-invariant
depth losses. It applies to reference outputs only, because the normal
sign convention is camera-specific.

\subsection{Albedo supervision}

For synthetic data with known albedo, use a logarithmic robust loss:
\[
  \Loss_\rho
  =
  \sum_v\sum_{\bm p\in\Omega_v}
  \left\|
    \log(\rho^v(\bm p)+\epsilon_\rho)
    -
    \log(\bar\rho^v(\bm p)+\epsilon_\rho)
  \right\|_1.
\]
Bounded albedo is enforced by the sigmoid output parameterization. In the initial
controlled setting, exposure, tone mapping, and renderer intensity settings
are fixed and recorded rather than freely learned per image. This reduces
confounding between reflectance, illumination, exposure, and geometry.

\subsection{Photometric shading loss with reliability weights}

For each ordered pair, compute the normal exclusively from the reference
view's camera-local pointmap and evaluate a single shading pathway:
\[
  \widehat I^a
  =
  \Sshade
  \left(
    X^{a,a},
    n_X(X^{a,a}),
    \rho^a,
    \mathcal Q_a,
    V
  \right).
\]
No pose is involved. Let $\psi(r)=\sqrt{r^2+\epsilon_\psi^2}$ be a robust penalty. For an RGB
residual, $\psi$ may be applied channel-wise and summed, or applied to a
specified robust RGB norm. Define
\[
  r^a(\bm p)
  =
  \psi\!\left(
    \widehat I^a(\bm p)-I_a(\bm p)
  \right).
\]

The shading loss is
\[
  \boxed{
  \Loss_{\mathrm{shade}}
  =
  \sum_{a}\sum_{\bm p\in\Omega_a}
  M_a(\bm p)
  \left[
    W^a(\bm p)r^a(\bm p)
    -
    \alpha_W\log(W^a(\bm p)+\epsilon_W)
  \right],
  }
\]
where the sum runs over the reference view of both orderings of each pair, and $M_a$ combines foreground, valid-depth,
finite-difference-validity, discontinuity, non-saturated, and non-dark masks. For a scalar residual $\psi$, the pointwise minimizer
of $W\psi-\alpha_W\log(W+\epsilon_W)$ over $0\leq W\leq1$ is
\[
  W^\star
  =
  \clip\!\left(\frac{\alpha_W}{\psi}-\epsilon_W,\,0,\,1\right).
\]
The logarithmic term discourages the reliability map from collapsing
toward zero. Reliability distributions and the fraction of active pixels
are monitored; the detached validity masks are not predicted by the network. The distribution of $W$ must be monitored during training to ensure
that the model does not suppress most photometric supervision.

\paragraph{What $W$ can and cannot mean.}
At the optimum, $W\psi\approx\alpha_W$, so the loss acts as a learned
iteratively reweighted robust penalty. $W$ cannot distinguish a model
violation (specular highlight, cast shadow) from a geometry or albedo
prediction error, since both raise the residual. It may down-weight pixels with large residuals regardless of whether the
cause is a model violation or a prediction error. The ablation against a
fixed robust (e.g.\ Cauchy) loss with no reliability head is essential; the
fixed robust loss is the training baseline.

The geometric normal is not detached by default: during training, the
photometric pathway differentiates through the finite-difference
normalization and into the pointmap. A stop-gradient variant may still
be evaluated as a training-stability ablation, but it is not a separate
regularizer.

\subsection{Deferred cross-view albedo consistency}

A future extension may use ground-truth or geometrically verified
correspondences to tie albedo predictions across views. This could supply a
multi-view photometric-stereo constraint, but it is excluded from the
initial objective and experiments to avoid coupling the first study to
correspondence quality and cross-view appearance consistency.
\section{Training-Time Visibility and Robustness}

The initial training objective uses the simple visibility approximation
\[
  V_k(\bm p)=1.
\]
For camera-attached lights close to the camera, cast shadows may be limited,
and the clamp $\max(0,\bm n^\top\bm\omega_k)$ handles attached-shadow
orientation but not occlusion by other scene geometry. Foreground, saturated,
near-black, invalid-depth, and discontinuity masks are applied where
appropriate. Robust photometric penalties and the learned reliability map
reduce the influence of some model violations, including cast shadows.
This is a deliberate simplification for the first training experiment, not
a claim that shadows are physically absent. Explicit differentiable
visibility and shadow tracing are deferred and are not part of either the
initial training objective or inference-time global alignment.

\section{Training Data and Protocol}

\subsection{Controlled synthetic data}

The initial dataset is generated in a metric simulator such as Gazebo. Scene
geometry, camera poses, and camera-relative light positions are specified
in meters. This establishes metric geometric ground truth, but does not by
itself provide radiometrically calibrated images. The initial study is
limited to Level 2: controlled illumination with recorded renderer settings.
The data-generation protocol should include:
\begin{itemize}
  \item diffuse-dominant materials with known RGB albedo metadata, using a
  consistent normalized reflectance convention;
  \item direct point lights, initially one effective light per image (or
  several simultaneous lights treated as one summed observation), with
  renderer intensity and attenuation settings recorded;
  \item camera-attached light positions expressed in camera coordinates,
  with randomized offsets/intensity settings in one configuration and a
  fixed-rig configuration in another;
  \item fixed exposure, tone mapping, color pipeline, and renderer version;
  if linear RGB output is available, retain it, otherwise document and
  consistently apply the renderer's fixed image-encoding conversion;
  \item known camera intrinsics, camera poses, metric depth/pointmaps,
  albedo metadata, and foreground masks;
  \item renderer normals may be exported for evaluation and diagnostics,
  but are not required as a training target.
\end{itemize}
The attenuation model in the analytic shading surrogate must be documented
and matched as closely as practical to the renderer configuration. It is
not described as physically calibrated inverse-square illumination. The
purpose is to test whether controlled illumination parameters offer
useful additional supervision for geometry. Training and test splits must
be separated by scene/object identity. Begin with simple diffuse scenes and
increase geometric clutter only after the basic experiment is stable.

\subsection{Mixing with general DUSt3R data}

The original DUSt3R training mixture may be retained to reduce forgetting
of general geometric capabilities. Samples without recorded illumination parameters
use the null-light condition and geometric objective only. Photometric
losses are applied only to samples with valid recorded light parameters and consistent renderer-image data. A frozen encoder, decoder LoRA, or gradual unfreezing
can be compared.

A crucial control is to fine-tune an otherwise equivalent DUSt3R model
on the same synthetic images without light conditioning and without
shading supervision. This separates the effect of controlled illumination
from the effect of additional synthetic training data.

\section{Experimental Plan and Evaluation}

\subsection{Progressive experiments}

All experiments are evaluated after the same standard geometric global
alignment procedure. The purpose is to measure whether training-time
photometric supervision improves the network's geometric predictions and,
in turn, the globally aligned reconstruction. No experiment introduces a
photometric objective into inference-time scene optimization.
\begin{enumerate}
  \item \textbf{E0 -- Geometric baselines:} evaluate the original DUSt3R
  checkpoint and a DUSt3R model fine-tuned on the same Gazebo images using
  geometric supervision only. The latter controls for the effect of
  additional synthetic training data. Report a newer pose-free backbone
  (e.g.\ MASt3R or VGGT) as an optional reference if resources permit.
  \item \textbf{E1 -- Oracle-albedo shading supervision:} fine-tune the
  geometric baseline with the derived-normal and shading losses, using
  ground-truth albedo only to evaluate the training-time shading operator.
  Ground-truth albedo is never supplied at inference. This isolates the
  potential geometric benefit of the shading gradient under favorable
  appearance information.
  \item \textbf{E1b -- Predicted-albedo supervision:} replace oracle albedo
  in the shading loss with the network's predicted albedo, supervised by
  synthetic albedo metadata. This tests the jointly learned appearance and
  geometry pathway under the intended inference setting.
  \item \textbf{E2 -- Reliability supervision:} compare the learned
  photometric-reliability head with a fixed robust penalty, keeping the
  geometric global-alignment procedure unchanged.
  \item \textbf{E3 -- Light-conditioned representation:} compare no explicit
  light conditioning, light tokens (self-attention only), and FiLM on both
  randomized-offset and fixed-rig data. All variants retain the same
  training-time geometric and photometric objectives where applicable.
\end{enumerate}
The first milestone is E0 followed by E1: establish a same-data
geometry-only control, then test whether the training-time shading signal
improves pointmaps before introducing predicted albedo, reliability, and
light-conditioning complexity.

\subsection{Metrics}

Report at least:
\begin{itemize}
  \item geometric accuracy, completeness, and symmetric Chamfer distance;
  \item depth error and depth-relative error;
  \item normal angular error, computed from the predicted pointmap against
  renderer normals when such evaluation data is available (not as a
  training requirement), reported alongside photometric error;
  
  \item relative camera pose metrics, such as rotation/translation
  accuracy and mean average accuracy where applicable;
  \item photometric reconstruction error in the fixed renderer-image encoding
  (and in linear RGB when directly available);
  \item sensitivity to perturbations in recorded light positions and renderer
  intensity/attenuation settings;
  \item the distribution of $W$ during training (to detect collapse or
  excessive suppression of photometric supervision).
\end{itemize}
Photometric error is a training diagnostic, not a substitute for geometric
evaluation. A model may reduce image residuals by changing albedo or
reliability without improving geometry. The primary evidence for the
research hypothesis must therefore come from pointmap and globally aligned
geometric metrics on held-out scenes.

\subsection{Required ablations}

\begin{center}
\begin{tabularx}{\textwidth}{>{\raggedright\arraybackslash}p{0.36\textwidth}X}
\toprule
Ablation & Question\\
\midrule
Original DUSt3R vs. geometry-only Gazebo fine-tuning &
How much improvement comes from additional synthetic geometric training?\\
Geometry-only vs. derived-normal supervision &
Does explicit supervision of pointmap derivatives improve surface orientation?\\
Oracle vs. predicted albedo in the training shading loss &
How much of the shading-supervision benefit survives albedo prediction?\\
Shading loss vs. no shading loss &
Does the image-formation constraint improve geometry beyond geometric supervision?\\
Learned reliability vs. fixed robust penalty &
Does the reliability head improve robustness to violations of the simplified model?\\
No light conditioning vs. light tokens vs. FiLM &
Where and how should illumination enter the network?\\
Randomized light offsets vs. fixed-rig illumination &
Does view-dependent light conditioning improve generalization?\\
Differentiable vs. stop-gradient derived normals &
Are gradients through pointmap finite differences stable and useful?\\
Accurate vs. perturbed light metadata &
How sensitive is training-time photometric supervision to metadata errors?\\
\bottomrule
\end{tabularx}
\end{center}

\subsection{Expected failure modes}

The following are hypotheses and anticipated limitations, not experimental
results:
\begin{itemize}
  \item Photometric supervision may improve surface orientation and fine-scale
  geometry more than coarse geometry.
  \item With a single effective light and no cross-view albedo tie, albedo
  and shading remain ambiguous and the normal is constrained only to a cone.
  \item Non-Lambertian materials, interreflection, and cast shadows may
  reduce the benefit of the v1 shading model.
  \item Noisy pointmap derivatives may make geometric normals unstable;
  validity-aware finite differences, controlled smoothing, derived-normal
  supervision, and detached-normal ablations help diagnose this issue.
  \item The shading loss may encourage the network to explain albedo or
  renderer mismatch through geometry; oracle/predicted-albedo ablations,
  geometric supervision, and held-out geometric metrics are needed to detect it.
  \item $W$ may down-weight pixels with large residuals caused by either
  model violations or prediction errors; it cannot distinguish the two.
  \item Incorrect light-position or renderer-parameter metadata may produce
  systematic geometric bias rather than merely higher photometric error.
\end{itemize}

\section{Inference Algorithm}

For an input image collection, the inference procedure deliberately retains
standard geometry-driven DUSt3R scene optimization:
\begin{enumerate}
  \item Provide the input images and, when the light-conditioned pathway is
  enabled, their recorded camera-relative light parameters to the trained
  LA-DUSt3R network.
  \item Run the network on both orderings of selected image pairs and retain
  the pairwise pointmaps and geometric-confidence predictions. Auxiliary
  albedo and photometric-reliability outputs may be returned for diagnostics,
  but are not used by scene optimization.
  \item Run the standard DUSt3R geometric global-alignment procedure using
  pairwise pointmaps and geometric confidence. Optimize the native geometric
  alignment variables to estimate camera poses, focal lengths, and a globally
  consistent reconstruction, preserving DUSt3R's original gauge and scale
  parameterization.
  \item Return the globally aligned 3D geometry, camera poses, and estimated
  focal lengths. Pointmap-derived normals and auxiliary albedo predictions
  may be computed or exported for evaluation, but they do not enter the
  global-alignment objective.
\end{enumerate}

No ground-truth poses or geometry are used at inference. The differentiable
shading operator, albedo loss, reliability-weighted photometric loss, and
any depth-residual optimization are training-only components and are not
evaluated during scene optimization. In particular, no photometric-family
constraint is added to the geometric global-alignment objective.

\section{Scope and Future Extensions}

The initial LA-DUSt3R study is intentionally limited to Level-2 controlled
synthetic illumination, camera-attached point lights with recorded
renderer-space parameters, and diffuse-dominant Lambertian-like materials.
The following extensions are deferred:
\begin{enumerate}
  \item differentiable visibility and screen-space shadow tracing;
  \item unknown-light or environment-map estimation;
  \item non-Lambertian reflectance, including specular BRDFs;
  \item indirect illumination and radiosity-inspired transport;
  \item cross-view albedo consistency and correspondence-driven
  multi-view photometric-stereo constraints;
  \item radiometrically calibrated illumination and recovery of absolute
  metric scale from photometry (metric scale is supplied by geometric
  supervision in the initial study);
  \item any post-alignment photometric optimization of depth, albedo, camera
  poses, or intrinsics; such refinement would require a separate study and
  is not part of the proposed inference pipeline;
  \item one-light-at-a-time multi-exposure capture, ambient-light handling
  for real flash images, and transfer to newer backbones.
\end{enumerate}
The differentiable light-transport and surfel literature provides
possible future directions, but its more elaborate transport solvers are
not required for the initial direct-illumination surrogate
\cite{gjerde2026point,jiang2025radiositygs}. A literature search on flash
photometric stereo and feed-forward relightable reconstruction should be
completed before claiming novelty.

\section{Summary}

LA-DUSt3R retains DUSt3R's pointmap-based geometric prior and augments the
training architecture with controlled-illumination conditioning,
diffuse-albedo and photometric-reliability prediction, pointmap-derived
normals with explicit derivative supervision, and a differentiable
Lambertian-like shading surrogate. The key consistency principle is to
evaluate the training-time shading loss in each image's own camera frame
using camera-local pointmaps from the swapped pair, without requiring camera
poses for the photometric pathway. Metric scale is supplied by metric
geometric supervision and is not recovered from illumination.

The separation between training and inference is deliberate. Photometric
constraints shape the learned geometric representation during training;
inference uses pairwise pointmaps and geometric confidence followed by the
standard DUSt3R geometric global alignment. No albedo, reliability, shading,
or depth-residual objective is introduced into scene optimization. The
initial experiments therefore test whether training-time photometric
supervision improves the globally aligned geometry, while preserving a
simple and well-understood inference procedure.

\appendix

\section{Revision notes for the revised v4.1}

\begin{itemize}
  \item Preserved the Level-2 controlled synthetic illumination scope and
  metric geometric supervision from the refined v4.
  \item Made the training/inference separation explicit: photometric losses
  are used only during network training.
  \item Removed the entire post-alignment photometric geometry-refinement
  stage, including depth residuals, albedo priors, photometric alignment
  objectives, and their optimization schedule.
  \item Restored standard DUSt3R geometric global alignment as the sole
  inference-time scene-optimization procedure, retaining its native gauge
  and scale parameterization.
  \item Clarified that inference-time global alignment uses pointmaps and
  geometric confidence; albedo and reliability are auxiliary outputs only.
  \item Reorganized experiments and ablations to measure the effect of
  training-time shading, albedo, reliability, and light conditioning after
  the same geometric global-alignment procedure.
  \item Retained metric scale from metric geometric supervision; metric-scale
  recovery from illumination and radiometric calibration remain deferred.
\end{itemize}

\section{Historical changes from v2 to v3 (superseded before revised v4.1)}

The following bullets document intermediate design history only. Any
post-alignment photometric optimization, depth-residual parameterization,
albedo optimization, or photometric alignment objective mentioned in these
notes is superseded and is not part of revised v4.1.

\begin{itemize}
  \item \textbf{Frames:} camera-local pointmaps come from the swapped pair
  (no pose); pose-based conversion $T_{2\leftarrow1}$ demoted to an ablation
  with the metric-scale caveat.
  \item \textbf{Scale:} added a scalar per-pair scale head, a global metric
  scale $s$ in alignment, and the prior
  $\Loss_{\mathrm{sc}}=\sum_e(\log s+\log\sigma_e-\log s_e)^2$ tying it to
  DUSt3R's $\sigma_e$ and gauge.
  \item \textbf{Identifiability:} new Sec.~\ref{sec:ident} (cone constraint,
  simultaneous lights sum, cross-view mechanism, scale--albedo degeneracy);
  optional cross-view albedo term restored (training and alignment).
  \item \textbf{Derivative supervision:} added $\Loss_{\mathrm{nrm}}$ using
  the same operator on the ground-truth pointmap; added a coarse-to-fine
  depth residual parameterization with penalty $\Loss_\delta$.
  \item \textbf{Masks:} discontinuity mask now uses 3D neighbor spacing
  versus pixel footprint, so grazing surfaces are retained.
  \item \textbf{Light conditioning:} default is light tokens in
  self-attention only (excluded from cross-attention keys); FiLM becomes an
  ablation; H3 qualified for rigid rigs.
  \item \textbf{Reliability:} documented what $W$ can and cannot mean; fixed
  the minimizer to include $\epsilon_W$; fixed robust loss as baseline.
  \item \textbf{Real data:} ambient-light and flash/no-flash note.
  \item \textbf{Experiments:} E1b (estimated albedo), newer backbones in E0,
  extended ablation table, $W$-distribution and normal-versus-photometric
  monitoring.
  \item \textbf{Notation and compile fixes:} single \texttt{xcolor} load;
  $\Loss_\rho/\lambda_\rho$ for training albedo, $\Loss_{\mathrm{ap}}/\lambda_{\mathrm{ap}}$
  for the alignment prior; $\Loss_{\mathrm{scale}}$ (training) distinguished from
  $\Loss_{\mathrm{sc}}$ (alignment); redundant inference step removed.
\end{itemize}

\begin{thebibliography}{9}

\bibitem{wang2024dust3r}
S. Wang, V. Leroy, Y. Cabon, B. Chidlovskii, and J. Revaud,
\newblock DUSt3R: Geometric 3D Vision Made Easy,
\newblock in \emph{Proceedings of the IEEE/CVF Conference on Computer
Vision and Pattern Recognition (CVPR)}, 2024.
\newblock arXiv:2312.14132.
\newblock \url{https://arxiv.org/abs/2312.14132}.
\newblock Code: \url{https://github.com/naver/dust3r}.

\bibitem{gjerde2026point}
M. K. Gjerde, J. B. Haurum, J. R. Frisvad, M. Worchel,
J. A. B{\ae}rentzen, and T. B. Moeslund,
\newblock Point-Based 3D Reconstruction from Sparse Views under Known
Illumination,
\newblock arXiv:2608.20000, 2026.
\newblock \url{https://arxiv.org/abs/2608.20000}.
\newblock Code: \url{https://github.com/M-Gjerde/DifferentiablePointRendering}.

\bibitem{jiang2025radiositygs}
K. Jiang, J.-M. Sun, Z. Li, D. Wang, T.-M. Li, and R. Ramamoorthi,
\newblock Differentiable Light Transport with Gaussian Surfels via
Adapted Radiosity for Efficient Relighting and Geometry Reconstruction,
\newblock \emph{ACM Transactions on Graphics}, SIGGRAPH Asia 2025.
\newblock arXiv:2509.18497.
\newblock \url{https://arxiv.org/abs/2509.18497}.
\newblock Code: \url{https://github.com/RaymondJiangkw/RadiosityGS}.

\end{thebibliography}

\end{document}
